\documentclass[10pt,a4paper]{amsart}
\usepackage[margin=19mm]{geometry}
\usepackage{amsmath,amssymb,mathtools}
\usepackage[scr=rsfs,cal=euler]{mathalfa}
\usepackage{xfrac}
\usepackage[unicode,hidelinks,bookmarks=false]{hyperref}
\newcommand{\ie}{i.e.\ }
\newcommand{\cf}{cf.\ }
\newcommand{\resp}{respectively\ }
\newcommand{\Subsection}{\subsection}
\providecommand{\nobreakdash}{\nobreak\hbox{-}}
\setlength{\emergencystretch}{2em}
\multlinegap0pt
\usepackage{xcolor}
\usepackage{amssymb}
\usepackage{mathtools}
\usepackage{tcolorbox}
\usepackage{algorithm}\usepackage{algpseudocode}
\newtheorem{corollary}{Corollary}
\newtheorem{lemma}{Lemma}
\usepackage{cleveref}
\crefname{corollary}{Corollary}{Corollarys}
\crefname{lemma}{Lemma}{Lemmas}
\newcommand{\defi}{\stackrel{\mathrm{\scriptscriptstyle def}}{=}}
\newcommand\linkofproof[1]{(\cref{#1})\newtarget{proof:#1}\ }
\newcommand\linktoproof[1]{{\hspace{-0.3em}\normalfont[{\hyperlink{proof:#1}{$\downarrow$}}]}}
\newcommand{\newtarget}[1]{\hypertarget{#1}{}}
\DeclareMathOperator{\sign}{sign}
\title{Curvature-Dependent Lower Bounds for Frank-Wolfe}
\author{Jannis Halbey, Christophe Roux, Sebastian Pokutta}
\date{Source: arXiv:2605.10595v2 (CC BY 4.0)}
\begin{document}
\maketitle

\noindent\emph{Source and licence.} Curvature-Dependent Lower Bounds for Frank-Wolfe. Version arXiv:2605.10595v2 (CC BY 4.0), distributed by arXiv under the Creative Commons Attribution 4.0 International licence, http://creativecommons.org/licenses/by/4.0/, as recorded in the arXiv licence metadata for this exact version and shown on the abstract page https://arxiv.org/abs/2605.10595v2. Full licence notice: \texttt{licenses/arxiv-2605.10595-CC-BY-4.0.txt}.\par\smallskip

\noindent\textit{Editorial scope.} Counted unit: the article's lemma lem:y-map (source0.tex lines 785--803). The article's complete original proof of it is delivered with it (source0.tex lines 804--1073, one \texttt{proof} environment of 270 lines). No mathematics is written, repaired or re-derived: the statement and the proof are the article's own lines, in the article's own notation, and no retained line is substituted or re-flowed.  Retained uncounted as the source-local context the proof needs: lines 709--717; lines 749--760.  Nothing was omitted from any retained span, so the package carries no omission record.  No retained line contains a citation, so the wrapper carries no bibliography and no bibliography entry.  No retained line contains a figure environment, an image-inclusion command or drawing code, so no image is delivered and the package declares no asset dependency.  Editorial formatting only, in addition: an AMS \texttt{amsart} preamble that restates the article's own declarations and macros used by the retained lines, namely the environments \texttt{corollary}, \texttt{lemma} on separate counters, together with the standard packages those lines need.  The retained environments print in this document as Corollary 1, Lemma 1 in order of appearance, whereas the article prints the same items with the numbering of its own full document.  The complete native source of the article as distributed in the arXiv e-print for this exact version is bundled unchanged as \texttt{sources/200308/source0.tex}.
    \begin{align}
        v_1 &= \frac{u^{q-1}}{M^{q-1}}, &
        v_2 &= -\frac{|w|^{q-1}}{M^{q-1}}\,\sign(w), \label{eq:lmo-uw} \\
        d_1 &= v_1-1+u, &
        d_2 &= v_2-w, \label{eq:dir-uw} \\
        \gamma &= \frac{M-u+u^2+w^2}{d_1^2+d_2^2}, \label{eq:gamma-uw} \\
        u' &= u-\gamma d_1, &
        w' &= w+\gamma d_2. \label{eq:update-uw}
    \end{align}

\begin{corollary}\label{cor:ratio-identity}
Assume $w\neq 0$ and define
$P \defi  |w|(1-v_1)+u|v_2|>0$. Then
\[
    u'=\frac{|d_2|\,P}{d_1^2+d_2^2},
    \qquad
    |w'|=\frac{|d_1|\,P}{d_1^2+d_2^2},
    \qquad
    \rho'=\frac{|w'|}{u'}=\frac{|d_1|}{|d_2|}.
\]
In particular, if $d_1\neq 0$ then $w'\neq 0$.
\end{corollary}

\begin{lemma}\label{lem:y-map}\linktoproof{lem:y-map}
    Fix $p\ge 3$ and let $q=p/(p-1)$, $\alpha=(p-1)/p$, and $\kappa=2\alpha$.
    Let $I\subset(0,p^{1/q})$ be compact.
    Then uniformly for $y\in I$, as $u\downarrow 0$,
    \begin{equation}\label{eq:y-map}
        \Phi(u,y)
        = F(y)+uG(y)+O(u^\kappa),
        \qquad
        F(y)=y^{-(q-1)}-\frac{y}{p},
        \qquad
        G(y)=\frac{y}{p}+\frac{p-1}{2p^2}\,y^{q+1}.
    \end{equation}
    Moreover, uniformly on $I$,
    \[
        \frac{\partial}{\partial y}\Phi(u,y)=F'(y)+uG'(y)+O(u^\kappa),
        \qquad
        \frac{\partial}{\partial u}\Phi(u,y)=G(y)+O(u^{\kappa-1}).
    \]
\end{lemma}

\begin{proof}\linkofproof{lem:y-map}
    The proof is divided into three parts.
    First, we are going to rewrite $\Phi(u,y)$ using some auxiliary functions that depend solely on $u$ and $y$ (or $z = y^q$).
    In the second part, we will expand each factor of the derived identity of $\Phi(u,y)$ around $u=0$.
    Finally, we show how these expansions extend to the partial derivatives of $\Phi(u,y)$.

    We begin by rewriting $\Phi(u,y)$ using \cref{cor:ratio-identity}
    \begin{equation}\label{eq:yprime-ratio}
        \Phi(u,y)= y'
        = \frac{\rho'}{(u')^\alpha}
        = \frac{|d_1|}{|d_2|(u')^\alpha}.
    \end{equation}
    In the following steps we will rewrite $d_1$, $d_2$ and $u'$ in terms of $u$ and $y$.
    First we define the auxiliary variable $A(u,y)\defi (1+y^q u)^{\frac1p}$.
    For any compact $I \subset (0,p^{1/q})$ and any $y \in I$, we have $z = y^q < p$.
    Using $\log(1+x) < x$ for $x > 0$, we have
    \[
        \log A(u,y) < \log (1+p u)^{\frac1p} = \frac1p \log (1+p u) < \frac1p (p u) = u,
    \]
    and therefore $\frac{1}{A(u,y)} > e^{-u}$.
    Next note that
    \begin{align}
        \label{eq:A-expand}
        A(u,y)&\defi (1+y^q u)^{\frac1p}
        = \left( 1 + \frac{|w|^q}{u^q}\right)^{\frac1p}
        =\frac{\left(u^q + |w|^q\right)^{\frac1p}}{u^{\frac qp}} \notag\\
        &= \left(\frac{M}{u}\right)^{\frac qp} = \left(\frac{M}{u}\right)^{q-1} = \frac{1}{v_1}.
    \end{align}
    Thus, by \eqref{eq:dir-uw} we have
    \begin{equation}
        \label{eq:d1-expand}
        d_1 = u + v_1 - 1 = u + \frac{1}{A} - 1 > u + e^{-u} - 1 > 0.
    \end{equation}
    For the second coordinate, we use $A = \left( \frac{M}{u}\right)^{q-1}$ from \eqref{eq:A-expand} to get
    \[
        |v_2| = \left(\frac{|w|}{M}\right)^{q-1} = \left(\frac{\rho u}{M}\right)^{q-1} = \frac{\rho^{q-1}}{A}
    \]
    and thus
    \begin{equation}
        \label{eq:d2-abs-expand}
        |d_2| = |v_2 - w|
        = |v_2|+|w|
        = \frac{\rho^{q-1}}{A} + y u^{1+\alpha}
        = \frac{(y u^\alpha)^{q-1}}{A} + y u^{1+\alpha}
        = \frac{y^{q-1}u^{\frac 1p}}{A} + y u^{1+\alpha}
    \end{equation}
    where we have used $\sign(v_2) = -\sign(w)$ and $\alpha(q-1)=\frac{q-1}{q}=\frac{1}{p}$.
    Next we define the auxiliary variables $E(u,y)$ and $\eta(u,y)$
    \[
        E(u,y)\defi y^{2-q}u^\kappa A(u,y),
        \qquad
        \eta(u,y)\defi \frac{\gamma d_1}{u}.
    \]
    Then we have
    \[
        |d_2|\,u^\alpha = \frac{y^{q-1}u}{A}\,(1+E),
        \qquad
        u'=u(1-\eta),
    \]
    so map $\Phi(u,y)$ admits the exact factorization
    \begin{equation}\label{eq:phi-factorization}
        \Phi(u,y) = \frac{d_1}{|d_2|(u')^\alpha}
        = \frac{d_1}{|d_2|u^\alpha(1-\eta)^\alpha}
        = \frac{A}{y^{q-1}}\frac{d_1}{u}(1+E)^{-1} (1-\eta)^{-\alpha},
    \end{equation}
    where we have dropped the absolute value for $d_1$ since $d_1 > 0$.

    Define
    \[
        B(u,y)\defi \frac{A}{y^{q-1}}\frac{d_1}{u},
        \qquad
        R(u,y)\defi (1+E)^{-1}(1-\eta)^{-\alpha},
    \]
    so that $\Phi=B\cdot R$ by \eqref{eq:phi-factorization}.
    We now expand the factors of $B$ and $R$ around $u=0$. By compactness of $I$,
    all estimates below are uniform for $y\in I$.
    After shrinking the $u \downarrow 0$ threshold, we may assume $zu=y^q u \leq \frac12$ for all $y \in I$.

    First we consider $A(u,y) = (1+zu)^{1/p}$.
    Since $A(0,y) = 1$ for any $y \in I$ and
    \begin{equation}
        \left.\frac{\partial }{\partial u} A(u,y)\right|_{u=0}
        = \left.\frac{1}{p} (1+zu)^{\frac 1p -1}z\right|_{u=0}
        = \frac{z}{p},
        \label{eq:A-deriv-1}
    \end{equation}
    we have
    \begin{equation}
        \label{eq:A-u-expand}
        A(u,y) = 1 + \frac{z}{p} u + \mathcal{O}(u^2).
    \end{equation}

    Next we compute the derivatives of $\frac{1}{A(u,y)}$ for $u=0$
    \begin{align}
        \left.\frac{\partial }{\partial u} \frac{1}{A(u,y)}\right|_{u=0}
        &= \left. -\frac{\frac{\partial }{\partial u} A(u,y)}{A(u,y)^2} \right|_{u=0}
        = -\frac{z}{p},
        \label{eq:invA-deriv-1}
        \\
        \left.\frac{\partial^2 }{\partial u^2} \frac{1}{A(u,y)}\right|_{u=0}
        &= \frac{2 \left(\frac{\partial }{\partial u} A(u,y) \right)^2 - A(u,y) \cdot \frac{\partial^2 }{\partial u^2} A(u,y)}{A(u,y)^3} \bigg|_{u=0} \notag \\
        &= 2 \left(\frac{z}{p}\right)^2 - 1 \cdot \frac{1-p}{p^2}z^2
        = \frac{p+1}{p^2}z^2,
        \label{eq:invA-deriv-2}
    \end{align}
    where we have used $A(0,y) = 1$ and \eqref{eq:A-deriv-1} as well as
    \begin{equation*}
        \left.\frac{\partial^2 }{\partial u^2} A(u,y)\right|_{u=0}
        = \left.\frac{1}{p}\left(\frac{1}{p}-1\right)(1+zu)^{\frac{1}{p}-2} z^2\right|_{u=0}
        = \frac{1-p}{p^2}z^2.
        \label{eq:A-deriv-2}
    \end{equation*}
    Thus, the expansion of $\frac{1}{A(u,y)}$ for $u=0$ yields
    \begin{align}
        \label{eq:d1-over-u-expand}
        \frac{d_1}{u}
        &= \frac{u + \frac{1}{A} - 1}{u}
        = \frac{u  + \left( 1 - \frac{z}{p} u + \frac12 \frac{p+1}{p^2}z^2 u^2 + O(u^3) \right) - 1}{u} \notag \\
        &= 1-\frac{z}{p} + \frac{p+1}{2p^2}z^2 u + O(u^2).
    \end{align}
    Combining \eqref{eq:A-u-expand} and \eqref{eq:d1-over-u-expand} yields
    \begin{align}
        B(u,y)
        &= \frac{A}{y^{q-1}}\frac{d_1}{u} \notag\\
        &= y^{-(q-1)} \left( 1 + \frac{z}{p} u + \mathcal{O}(u^2) \right) \left( 1 - \frac{z}{p} + \frac{p+1}{2p^2}z^2 u + O(u^2) \right) \notag\\
        &= F(y) + uG(y) + O(u^2). \label{eq:B-u-expand}
    \end{align}
    Next we consider the terms of $R(u,y)$ starting with $(1+E)^{-1}$.
    Using the expansion $(1+x)^{-1} = 1 - x + O(x^2)$ for $x \to 0$ and
    \[
        E(u,y) = y^{2-q}u^\kappa A(u,y) = y^{2-q}u^\kappa \left( 1 + \frac{z}{p} u + \mathcal{O}(u^2) \right) = \mathcal{O}(u^\kappa),
    \]
    we get
    \begin{equation}
      \label{eq:1+E-inv-expand}
        (1+E)^{-1} = 1 + O(u^\kappa).
    \end{equation}
    In the next step we expand $(1-\eta)^{-\alpha}$.
    For that we rewrite $\eta$ using \eqref{eq:gamma-uw},
    \[
        \eta = \frac{\gamma d_1}{u} = \frac{d_1}{u} \cdot \frac{M - u + h}{d_1^2+d_2^2}.
    \]
    First we see that
    \[
        u(1+zu)^\alpha = u \left( 1 + \frac{|w|^q}{u^q}\right)^{\frac 1q} = (u^q + |w|^q)^{\frac 1q} = M,
    \]
    and
    \[
        h = u^2 + w^2 = u^2 + y^2 u^{2(1+\alpha)} = u^2 + y^2 u^{2+\kappa}.
    \]
    Expansion yields $(1+zu)^\alpha = 1 + \alpha z u + O(u^2)$ and thus we get
    \begin{align}
        \label{eq:M-u+h-expand}
        M-u+h
        &= u(1+zu)^\alpha - u + u^2 + y^2 u^{2+\kappa} \notag\\
        &= u(1+\alpha z u + \mathcal{O}(u^2)) - u + u^2 + y^2 u^{2+\kappa} = \mathcal{O}(u^2).
    \end{align}
    Furthermore, by \eqref{eq:d1-over-u-expand} we have $d_1 = \mathcal{O}(u)$
    and \eqref{eq:d2-abs-expand} and \eqref{eq:A-u-expand} imply $|d_2| = \Theta(u^{1/p})$.
    Therefore, together with \eqref{eq:M-u+h-expand},
    \begin{equation}
        \label{eq:gamma-expand}
        \gamma = \frac{M - u + h}{d_1^2+d_2^2} = \frac{\mathcal{O}(u^2)}{\Theta(u^{2/p})} = \mathcal{O}(u^{2- 2/p}) = \mathcal{O}(u^\kappa).
    \end{equation}
    Together with \eqref{eq:d1-over-u-expand} this yields
    \[
        \eta = \frac{d_1}{u} \cdot \gamma = \mathcal{O}(1) \cdot \mathcal{O}(u^\kappa) = \mathcal{O}(u^\kappa).
    \]
    With $(1-x)^{-\alpha} = 1 + \alpha x + O(x^2)$ for $x \to 0$ we get
    \begin{equation}
      \label{eq:1-eta-alpha-expand}
        (1-\eta)^{-\alpha} = 1 + \alpha \eta + O(\eta^2) = 1 + \mathcal{O}(u^\kappa).
    \end{equation}
    Therefore we have
    \begin{equation}
        \label{eq:R-u-expand}
        R(u,y)= (1+E)^{-1}(1-\eta)^{-\alpha} = (1+O(u^\kappa))(1+O(u^\kappa)) = 1+O(u^\kappa).
    \end{equation}
    Combining with \eqref{eq:B-u-expand} yields
    \begin{align}
        \label{eq:phi-u-expand}
        \Phi(u,y) &= B(u,y)R(u,y) = (F(y) + uG(y) + O(u^2))(1+O(u^\kappa)) \notag\\
        &= F(y) + uG(y) + O(u^2) + O(u^\kappa) = F(y) + uG(y) + O(u^\kappa),
    \end{align}
    as $\kappa = 2\cdot \frac{p-1}{p} < 2$.

    In the last step, we derive expansions for the partial derivatives of $\Phi(u,y)$ starting with $B$.
    First note that $A(u,y)=(1+y^q u)^{1/p}$ is smooth and $A(u,y) > 0$ for all $(u,y) \in [0,u_0]\times I$
    and thus $d_1 = u + \frac 1A - 1$ is smooth.
    Furthermore, since $d_1(0,y) = 0$ for all $y \in I$, the quotient $\frac{d_1}{u}$ is also smooth and finally $B(u,y)=y^{-(q-1)}A(u,y)\frac{d_1}{u}$ is smooth.
    Since $[0,u_0]\times I$ is compact, all second partial derivatives of $B(u,y)=F(y)+uG(y)+O(u^2)$ are bounded there
     and thus the $O(u^2)$ remainder stays bounded after differentiation.
    This yields
    \[
        \frac{\partial}{\partial y} B(u,y)=F'(y)+uG'(y)+O(u^2),\qquad
        \frac{\partial}{\partial u} B(u,y)=G(y)+O(u).
    \]
    uniformly for $y \in I$.

    Next we estimate derivatives of $R=(1+E)^{-1}(1-\eta)^{-\alpha}$.
    Since
    \[
        E(u,y)=y^{2-q}u^\kappa A(u,y),
    \]
    with $A=O(1)$, $\frac{\partial}{\partial y} A=O(u)$ and $\frac{\partial}{\partial u} A=O(1)$ uniformly on $I$, product
    rule gives
    \[
      \frac{\partial}{\partial y} E(u,y)=O(u^\kappa),\qquad
      \frac{\partial}{\partial u} E(u,y)=O(u^{\kappa-1}).
    \]
    Since $A >0$, we have $E(u,y) \geq 0$ and thus $(1+E)^{-2} = \mathcal{O}(1)$.
    The chain rule yields then
    \[
        \frac{\partial}{\partial y} (1+E)^{-1}=-(1+E)^{-2}\frac{\partial}{\partial y} E=O(u^\kappa),\qquad
        \frac{\partial}{\partial u} (1+E)^{-1}=-(1+E)^{-2}\frac{\partial}{\partial u} E=O(u^{\kappa-1}).
    \]
    Next we write $\eta=\frac{d_1}{u}\gamma$ with
    \[
        \gamma=\frac{M-u+h}{d_1^2+d_2^2},
    \]
    Since $M$ and $h$ are both twice continuously differentiable with $M - u + h  = \mathcal{O}(u^2)$, we get
    \[
        \frac{\partial}{\partial y} (M-u+h)=O(u^2),\qquad
        \frac{\partial}{\partial u} (M-u+h)=O(u).
    \]
    Further, $d_1=O(u)$ and \eqref{eq:d2-abs-expand} imply
    \[
        d_1^2+d_2^2=\Theta(u^{2/p}),\quad
        \frac{\partial}{\partial y}(d_1^2+d_2^2)=O(u^{2/p}),\quad
        \frac{\partial}{\partial u}(d_1^2+d_2^2)=O(u^{2/p-1}).
    \]
    Applying the quotient rule to $\gamma$ gives
    \[
        \frac{\partial}{\partial y} \gamma=O(u^\kappa),\qquad
        \frac{\partial}{\partial u} \gamma=O(u^{\kappa-1}),
    \]
    since $\kappa=2-2/p$.
    Furthermore, \eqref{eq:d1-over-u-expand} and the smoothness of $\frac{d_1}{u}$ yield
    \[
        \frac{\partial}{\partial y} \left(\frac{d_1}{u}\right)=O(1),\qquad
        \frac{\partial}{\partial u} \left(\frac{d_1}{u}\right)=O(1).
    \]
    Therefore, by product rule for $\eta=\frac{d_1}{u}\gamma$ we get
    \[
        \frac{\partial}{\partial y} \eta=O(u^\kappa),\qquad
        \frac{\partial}{\partial u} \eta=O(u^{\kappa-1}).
    \]
    For $u_0$ small enough we have $|\eta|\le \frac12$, hence
    $(1-\eta)^{-\alpha-1}=O(1)$ uniformly, and chain rule yields
    \[
        \frac{\partial}{\partial y} (1-\eta)^{-\alpha}=O(u^\kappa),\qquad
        \frac{\partial}{\partial u} (1-\eta)^{-\alpha}=O(u^{\kappa-1}).
    \]
    Combining with the bounds for $(1+E)^{-1}$ and its derivatives,
    \[
        \frac{\partial}{\partial y} R(u,y)=O(u^\kappa),\qquad
        \frac{\partial}{\partial u} R(u,y)=O(u^{\kappa-1}).
    \]
    Finally, from $\Phi=B\cdot R$, we get 
    \[
        \frac{\partial}{\partial y} \Phi(u,y)=(\frac{\partial}{\partial y} B)R+B\,\frac{\partial}{\partial y} R,\qquad
        \frac{\partial}{\partial u} \Phi(u,y)=(\frac{\partial}{\partial u} B)R+B\,\frac{\partial}{\partial u} R.
    \]
    Since $B=F+uG+O(u^2)=O(1)$ and all bounds are uniform on $I$, we get
    \[
        \frac{\partial}{\partial y} \Phi(u,y)=F'(y)+uG'(y)+O(u^\kappa),\qquad
        \frac{\partial}{\partial u} \Phi(u,y)=G(y)+O(u^{\kappa-1}).
    \]
    This completes the proof.
\end{proof}

\end{document}
